\section{Upstream Code and Corrections}\label{app:upstream}

\paragraph{Sources.} The localization filters, the resampling schemes and the rules for when to resample come from the
particle filter tutorial of Elfring et al.\ \cite{elfring2021particle} (\url{https://github.com/jelfring/particle-filter-tutorial},
commit \texttt{e6014b7}, MIT licence). FastSLAM 1.0 and 2.0 come from PythonRobotics \cite{sakai2018pythonrobotics} (\url{https://github.com/AtsushiSakai/PythonRobotics},
\texttt{SLAM/FastSLAM1} and \texttt{SLAM/FastSLAM2}, commit \texttt{1fe4fb9}, MIT licence). The upstream code is
downloaded at these commits and never modified. The resampling schemes and rules are called unchanged; the models,
proposals and FastSLAM updates are copied, vectorized over particles, and changed to work with log weights.

\paragraph{Checking.} The corrections to the models and proposals can be switched off to restore the upstream behaviour.
With them switched off, the copies were compared with the upstream functions on the same inputs and random numbers. The localization likelihood and the noise-free
motion model match exactly, and one full extended Kalman particle filter update matches in poses, covariances and weights
when both use the Gaussian mean instead of a random draw. The FastSLAM prediction matches exactly with the same seed, and
the FastSLAM 1.0 and 2.0 observation updates match in weights, landmark means and covariances over two steps. These checks
are part of the test suite.

\paragraph{Corrections in the localization code.}
\begin{enumerate}
  \item The likelihood compares the expected and measured bearing without wrapping the difference, so a small error
        across $\pm\pi$ looks like an error of almost $2\pi$ and the particle gets almost no weight. The difference is
        wrapped to $[-\pi, \pi)$.
  \item The likelihood is the exponent of a Gaussian without its normalizing constant. The constant is the same for
        every particle, so it cancels when the weights are normalized and no result of
        \crefrange{sec:e01}{sec:e04} depends on it, but it is exactly the term that penalizes a wider assumed noise.
        Without it the marginal likelihood \eqref{eq:evidence} grows without bound as the assumed measurement noise
        grows, and no parameter can be estimated from it. The full log density is used.
  \item Extended Kalman particle filter: the covariance prediction multiplies element by element instead of computing
        the matrix product $F P F^\top$; the noise covariances are built from standard deviations instead of variances;
        the bearing innovation is not wrapped; the prior is centred on the updated state instead of the prediction,
        because the predicted state and the updated state are the same list object; and the bearing Jacobian divides
        by zero when the particle and the landmark have the same $x$ coordinate. All five are corrected.
  \item Extended Kalman particle filter: positions wrap around at the edges of the world, but the prior and proposal
        densities use raw position differences. A pose sampled just across an edge then appears a world-width away from
        its prediction and can receive an arbitrarily large weight; from a uniform start this makes all particles
        collapse onto one wrong pose more often the more particles there are. The differences are wrapped to the world.
        A test checks that the mean weight from a uniform start equals the mean likelihood under the prior, as importance
        sampling requires; without the correction the two differ by a factor of more than $e^{10^4}$.
  \item The extended Kalman particle filter hard-codes multinomial resampling at every step, and the auxiliary particle
        filter hard-codes multinomial sampling in its first stage. Resampling is left to the configured scheme and rule,
        and the auxiliary filter's first-stage selection uses the configured scheme.
  \item The systematic and stratified search loops can step past the last particle when rounding leaves the cumulative
        weight slightly below one. The last weight is raised by that shortfall before calling them.
\end{enumerate}

\paragraph{Corrections in the FastSLAM code.}
\begin{enumerate}
  \item The FastSLAM 2.0 proposal moves each particle to the mean of its proposal Gaussian but never draws a sample from
        it, never corrects the weight for the proposal, and keeps the proposal covariance from one step to the next. This
        is closer to a Kalman update per particle than to FastSLAM 2.0, and comparing it with FastSLAM 1.0 would not compare
        a motion-model proposal with a measurement-informed one. The corrected proposal follows Montemerlo et al.\
        \cite{montemerlo2003fastslam2}: it predicts the pose, updates the pose Gaussian with each observation of an already
        mapped landmark while weighting by the predictive likelihood, samples the pose, and then updates the landmark filters.
  \item Resampling copies the list of particle objects but not the particles, so resampled particles share their state and
        their landmark arrays. Particles are stored as arrays and copied on resampling.
  \item A landmark counts as new when its estimated $x$ coordinate is within 0.01 of zero, so a landmark estimated near
        $x = 0$ is initialized again. A flag per particle and landmark records which landmarks have been seen.
  \item The pose estimate averages headings as plain numbers, which fails near $\pm\pi$. A circular mean is used.
  \item Resampling is hard-coded (low-variance resampling when $\mathrm{ESS} < N / 1.5$). The configured scheme and rule
        are used, with this rule as the default.
\end{enumerate}

\paragraph{Effect.} With the upstream code, the comparison of FastSLAM 1.0 and 2.0 in \cref{sec:e03} would have compared
two variants of a motion-model proposal. Without the normalizing constant the sweeps of \cref{sec:e11} rise
monotonically with the assumed measurement noise and have no maximum. Without the correction at the world's edges, the extended Kalman particle filter
becomes less accurate as particles are added. The results in \cref{sec:experiments} use the corrected code throughout.
